\begin{itemframe}{مبانی موازی‌سازی انشعاب-الحاق}
\itm
بحث برنامه‌نویسی موازی را با مسئله‌ی محاسبه‌ی اعداد فیبوناچی به‌صورت بازگشتی آغاز می‌کنیم.
\itm
ابتدا به یک محاسبه‌ی ساده‌ی دنباله فیبوناچی نگاه می‌کنیم که هرچند ناکارآمد است، اما نحوه‌ی بیان موازی‌سازی در شبه‌کد را نشان می‌دهد.
\end{itemframe}


\begin{itemframe}{مبانی موازی‌سازی انشعاب-الحاق}
\itm
الگوریتم
\texttt{FIB}
عدد فیبوناچی $n$ام را به‌صورت بازگشتی محاسبه می‌کند.
\begin{algorithm}[H]\alglr
\caption{FIB}
\begin{algorithmic}[1]
\Func{FIB}{$n$}
\If{$n \leq 1$}
    \State \Return $n$
\Else
    \State $x \gets \text{FIB}(n-1)$
    \State $y \gets \text{FIB}(n-2)$
    \State \Return $x + y$
\EndIf
\end{algorithmic}
\end{algorithm}

\end{itemframe}


\begin{itemframe}{مبانی موازی‌سازی انشعاب-الحاق}
\itm
برای تحلیل این الگوریتم، فرض کنید $T(n)$ زمان اجرای FIB(n) را نشان دهد.
 FIB(n)
شامل دو فراخوانی بازگشتی به‌علاوه‌ی مقداری کار ثابت است:
$$
T(n) = T(n-1) + T(n-2) + \Theta(1)
$$
\itm
جواب این رابطه‌ی بازگشتی با استفاده از روش جایگزینی
\fn{substitution method}
به دست می‌آید که نتیجه می‌دهد:
$$
T(n) = \Theta(\phi^n)
$$
که در آن
$$
\phi = \frac{1+\sqrt{5}}{2}
$$
\end{itemframe}


\begin{itemframe}{مبانی موازی‌سازی انشعاب-الحاق}
\itm
دیدم که زمان اجرای این روش نمایی است، بنابراین بسیار کند است. بیایید ببینیم چرا این الگوریتم ناکارآمد است. شکل زیر درخت فراخوانی‌های بازگشتی را نشان می‌دهد که توسط الگوریتم FIB ایجاد می‌شوند.
\pic[.8]{1}
\end{itemframe}


\begin{itemframe}{مبانی موازی‌سازی انشعاب-الحاق}
\itm
گرچه رویه‌ی FIB روش ناکارآمدی است، می‌تواند به ما کمک کند تا با مفاهیم موازی‌سازی آشنا شویم.
\itm
شاید ابتدایی‌ترین مفهوم این است که اگر دو وظیفه‌ی موازی روی داده‌های کاملاً مستقل عمل کنند
آنگاه تفاوتی نمی‌کند که این دو همزمان اجرا شوند یا به صورت سریالی و یکی پس از دیگری.
در هر دو صورت نتیجه یکسان است.
\itm
به‌عنوان مثال، در الگوریتم FIB دو فراخوانی بازگشتی
$FIB(n-1)$
و
$FIB(n-2)$
می‌توانند با ایمنی کامل به صورت موازی اجرا شوند، زیرا هیچ‌یک اثری بر محاسبات دیگری ندارد.
\end{itemframe}


\begin{itemframe-s}{مبانی موازی‌سازی انشعاب-الحاق}{کلیدواژه‌های موازی‌سازی}
\itm
رویه‌ی P-FIB، اعداد فیبوناچی را محاسبه می‌کند، اما در شبه‌کد خود از کلیدواژه‌های spawn و sync برای نمایش موازی‌بودن استفاده می‌کند.
\begin{algorithm}[H]\alglr
\caption{P-FIB}
\begin{algorithmic}[1]
\Func{P-FIB}{$n$}
\If{$n \leq 1$}
    \State \Return $n$
\Else
    \State $x \gets \textbf{spawn }$ P-FIB($n-1$) \Comment{don’t wait for subroutine to return}
    \State $y \gets$ P-FIB($n-2$) \Comment{in parallel with spawned subroutine}
    \State \textbf{sync} \Comment{wait for spawned subroutine to ûnish}
    \State \Return $x+y$
\EndIf
\end{algorithmic}
\end{algorithm}
\end{itemframe-s}


\begin{itemframe-s}{مبانی موازی‌سازی انشعاب-الحاق}{کلیدواژه‌های موازی‌سازی}
\itm
اگر کلیدواژه‌های spawn و sync از رویه‌ی P-FIB حذف شوند، متن شبه‌کد حاصل دقیقاً مشابه رویه‌ی FIB خواهد بود.
\itm
«تصویر سریالی»
\fn{serial projection}
الگوریتم موازی، یک الگوریتم سریالی‌ است که با حذف دستورات موازی‌سازی به‌دست می‌آید.
\itm
اگر یک الگوریتم موازی داشته باشیم آنگاه، تصویر سریالی آن همواره یک شبه‌کد سریالی برای حل همان مسئله است.
\end{itemframe-s}


\begin{itemframe-s}{مبانی موازی‌سازی انشعاب-الحاق}{گراف توازی}
\itm
زمانی که کلیدواژه‌ی spawn قبل از یک فراخوانی رویه بیاید، می‌گوییم «فراخوانی موازی»
\fn{Spawning}
رخ می‌دهد؛ مانند خط ۳ در
\texttt{P-FIB}
\itm
رویه‌ای که دستور spawn را اجرا می‌کند (والد)، می‌تواند به‌طور موازی با زیررویه‌ی ایجادشده (فرزند) ادامه یابد. درحالی که در حالت سریالی والد باید منتظر بماند تا اجرای فرزند تمام شود.
\itm
در این مثال، در حالی که فرزند ایجادشده در حال محاسبه‌ی
$P-FIB(n - 1)$
است، والد می‌تواند در خط ۴ محاسبه‌ی
$P-FIB(n - 2)$
را موازی با فرزند اجرا کند.
\end{itemframe-s}


\begin{itemframe-s}{مبانی موازی‌سازی انشعاب-الحاق}{گراف توازی}
\itm
از آنجا که رویه‌ی P-FIB بازگشتی است، این دو فراخوانی خودشان نیز فرزندهای موازی ایجاد می‌کنند، و فرزندانشان نیز همین کار را انجام می‌دهند؛ بنابراین یک درخت بسیار بزرگ شکل می‌گیرد که همگی به‌طور موازی اجرا می‌شوند.
\itm
البته کلیدواژه‌ی spawn به این معنی نیست که یک رویه الزاماً باید موازی با فرزندانش اجرا شود، تنها می‌گوید یک فراخوانی می‌تواند موازی اجرا شود.
\itm
به طور کلی کلیدواژه‌های موازی‌سازی، موازی‌بودن محاسبات را به طور «منطقی»
\fn{logical parallelism}
بیان می‌کنند.
یعنی نشان می‌دهند که کدام بخش‌ها می‌توانند موازی اجرا شوند.

\end{itemframe-s}


\begin{itemframe-s}{مبانی موازی‌سازی انشعاب-الحاق}{گراف توازی}
\itm
در زمان اجرا، واحدی به نام زمان‌بند
\fn{scheduler}
تعیین می‌کند کدام زیر محاسبات
\fn{subcomputations}
واقعاً موازی اجرا شوند.
\itm
همچنین یک رویه نمی‌تواند از مقادیر بازگشتی فرزندان خود استفاده کند، مگر پس از اجرای یک دستور sync (مانند خط ۵).
\itm
کلیدواژه‌ی sync بیان می‌کند که رویه باید منتظر بماند تا همه‌ی فرزندان آن پایان یابند، سپس به اجرا ادامه دهد.
\itm
علاوه بر همگام‌سازی صریح با استفاده از دستور sync این فرض ساده‌کننده را نیز می‌پذیریم که هر رویه پیش از دستور return به‌طور ضمنی یک sync اجرا می‌کند؛ این امر تضمین می‌کند که همه‌ی فرزندان پیش از پایان والد به اتمام برسند.
\end{itemframe-s}


\begin{itemframe-s}{مبانی موازی‌سازی انشعاب-الحاق}{گراف توازی}
\itm
برای درک اجرای یک محاسبه‌ی موازی می‌توان آن را به‌صورت یک گراف جهت‌دار بدون دور نشان داد. این گراف یک که یک گراف توازی
\fn{parallel trace}
نامیده می‌شود.
\itm
به طوری که رأس‌های این گراف نشان‌دهنده‌ی دستورالعمل‌های اجراشده، و یال‌های آن وابستگی بین دستورالعمل‌ها را نمایش می‌دهند، به‌طوری که اگر
$(u, v) \in E$
باشد، به این معناست که دستور $u$ باید پیش از دستور $v$ اجرا شود.
\itm
اما نمایش هر دستور به‌عنوان یک رأس در گراف توازی، گراف را بی دلیل شلوغ می‌کند.
بنابراین اگر دنباله‌ایی از دستورها شامل هیچ کلیدواژه موازی‌سازی یا دستور کنترل توابع نباشد (به طور دقیق spawn، sync، فراخوانی و بازگشت توابع)، کل این زنجیره را در قالب یک بلوک
\fn{strand}
گروه‌بندی می‌کنیم.
\end{itemframe-s}


\begin{itemframe-s}{مبانی موازی‌سازی انشعاب-الحاق}{گراف توازی}
\itm
به‌عنوان مثال، شکل زیر یک گراف توازی که از اجرای
$P-FIB(4)$
حاصل می‌شود را نشان می‌دهد.
\pic[.6]{2}
\end{itemframe-s}


\begin{itemframe-s}{مبانی موازی‌سازی انشعاب-الحاق}{گراف توازی}
\itm
در شکل، هر دایره یک بلوک‌ی محاسباتی را نشان می‌دهد؛ دایره‌های آبی شامل دستورات تابع FIB از اول تا فراخوانی P-FIB(n-1) هستند.
\itm
دایره‌های نارنجی از فراخوانی P-FIB(n-2) شروع و با دستور sync پایان می‌یابند.
و دایره‌های سفید شامل دستورات بعد از پس از sync  هستند.
 بلوک‌هایی که متعلق به یک تابع هستند گروه‌بندی می‌شوند؛ رنگ آبی برای توابع فراخوانی‌شده با spawn و رنگ قهوه‌ای برای توابعی که فراخوانی معمولی شده اند.
\itm
یال‌های نشان‌داده‌شده با رنگ آبی مسیر بحرانی را نشان می‌دهند که در آینده بیشتر با آن آشنا خواهیم شد.
\end{itemframe-s}


\begin{itemframe-s}{مبانی موازی‌سازی انشعاب-الحاق}{گراف توازی}
\itm
همانطور که دیدیم بلوک‌ها شامل دستورهای کنترل توازی یا کنترل توابع نمی‌باشند؛ بلکه این وابستگی‌های به‌صورت یال نمایش داده شوند.
\itm
هنگامی که یک والد یک فرزند را فراخوانی می‌کند، گراف توازی شامل یالی
از بلوک‌ی که فراخوانی را اجرا می‌کند در والد به نخستین بلوک‌ی فرزند ایجادشده می‌کنیم.

 این حالت در شکل، توسط یک یال از بلوک‌ی نارنجی در
$P-FIB(4)$
به بلوکی آبی در
$P-FIB(2)$
نشان داده شده است.
\itm
زمانی که آخرین بلوک در فرزند بازمی‌گردد، یالی از آخرین بلوک تابع فرزند به بلوک بعدی در والد (بلوک بعد از فراخوانی کننده) اضافه می‌کنیم.
نمونه‌ی آن یال از بلوک‌ی سفید در
$P-FIB(2)$
به بلوک‌ی سفید در
$P-FIB(4)$
است.
\end{itemframe-s}


\begin{itemframe-s}{مبانی موازی‌سازی انشعاب-الحاق}{گراف توازی}
\itm
می‌توانید به یال‌هایی که در بالا توصیف شد در شکل دقت کنید:
\pic[.6]{2}
\end{itemframe-s}


\begin{itemframe-s}{مبانی موازی‌سازی انشعاب-الحاق}{گراف توازی}
\itm
چیزی که در بالا توصیف شد فراخوانی عادی بود. گفتیم زمانی که قبل از فراخوانی کلمه spawn بیاید تبدیل به فراخوانی موازی می‌شود. بیایید فراخوانی موازی را نیز در شکل بررسی کنیم.
\itm
مانند فراخوانی عادی، در اینجا هم از والد به فرزند یال وجود دارد، اما یک یال دیگر نیز وجود دارد
که نشان می‌دهد در فرزند، اجرا می‌تواند به بلوک بعدی انتقال یابد در حالی که بلوک فراخوانی کننده در حال اجراست.
\itm
به‌عنوان مثال، یال از بلوک‌ی آبی در
$P-FIB(4)$
به بلوک‌ی نارنجی در
$P-FIB(4)$
نمونه‌ای از چنین یالی است.
\itm
همچنین، مانند فراخوانی عادی، از آخرین بلوک فرزند یک یال خارج می‌شود؛ که در فراخوانی عادی به بلوک بعد از فراخوانی کننده می‌رود.
اما در فراخوانی موازی به بلوک بعد از دستور sync در والد می رود.
\end{itemframe-s}


\begin{itemframe-s}{مبانی موازی‌سازی انشعاب-الحاق}{گراف توازی}
\itm
اگر در گراف توازی مسیری جهت‌دار از $u$ به $v$ داشته باشد، می‌گوییم این دو بلوک «سریِ منطقی» هستند.
\fn{logically in series}
اگر هیچ مسیری نه از $u$ به $v$ و نه از $v$ به $u$ وجود نداشته باشد، این دو بلوک «موازیِ منطقی»
\fn{logically in parallel}
هستند.
\itm
در تحلیل‌ها فرض می‌کنیم که کامپیوتر ما شامل مجموعه‌ای از پردازنده‌ها و یک حافظه‌ی مشترک «سازگار ترتیبی»
\fn{sequentially consistent}
است.
\itm
برای اینکه سازگاری ترتیبی را بهتر درک کنیم باید کمی در مورد دستورات سطح پایین (اسمبلی) کار با حافظه بیشتر بدانیم.
\end{itemframe-s}


\begin{itemframe-s}{مبانی موازی‌سازی انشعاب-الحاق}{گراف توازی}
\itm
دستورهای دسترسی به حافظه به دو صورت‌اند:\\
load
:کپی داده از حافظه به ثبات پردازنده \\
store
:کپی داده از ثبات پردازنده به حافظه
\itm
برای نمونه، دستور زیر شامل چندین load و store است:
$$x = y + z$$

\itm
در یک رایانه‌ی موازی، ممکن است چند پردازنده هم‌زمان به حافظه دسترسی داشته باشند. سازگاری ترتیبی سلامت داده‌ها را تضمین می‌کند.
\end{itemframe-s}


\begin{itemframe-s}{مبانی موازی‌سازی انشعاب-الحاق}{گراف توازی}
\itm
درواقع رفتار حافظه مشترک طوری است که گویی در هر لحظه تنها یک دستور از یکی از پردازنده‌ها اجرا شده است.
در حالی که ممکند است چندین انتقال داده به‌طور فیزیکی هم‌زمان رخ دهند.
\itm
محاسبات وظیفه-موازی، توسط سیستم زمان‌ روی پردازنده‌ها مدیریت می‌شوند. در این حالت ترتیب اجرای دستورات روی حافظه‌ مطابق یک مرتب‌سازی توپولوژیکی از گراف توازی است.
\itm
این بدان معناست که می‌توان اجرای واقعی را به‌صورت یک توالی خطی از دستورها تصور کرد.
اگرچه زمان‌بند، می‌تواند ترتیب اجرا را بین دفعات مختلف اجرا تغییر دهد، نتیجهٔ محاسبه همواره معادل اجرای سریالی است.

\end{itemframe-s}


\begin{itemframe-s}{مبانی موازی‌سازی انشعاب-الحاق}{گراف توازی}
\itm
علاوه بر فرض‌های قبلی در مورد رایانه مورد استفاده، دو فرض دیگر هم اضافه می‌کنم:
\item[1]
همه‌ی پردازنده‌ها توان محاسباتی برابر دارند.
\item[2]
هزینه‌ی زمان‌بندی نادیده گرفته می‌شود.
\itm
هرچند این فرض دوم خوش‌بینانه به نظر می‌رسد، در عمل با یک پیاده‌سازی مناسب، سربار زمان‌بندی ناچیز است.

\end{itemframe-s}


\begin{itemframe-s}{مبانی موازی‌سازی انشعاب-الحاق}{معیارهای کارایی}
\itm
می‌توانیم کارایی یک الگوریتم وظیفه-موازی را با استفاده از تحلیل کار-گستره
\fn{work/span analysis}
بسنجیم؛ تحلیلی که بر اساس دو معیار کار
\fn{work}
و گستره
\fn{span}
تعریف می‌شود.
\itm
«کار» کل زمانی است که طول می‌کشد تا کل محاسبه بر روی
\textbf{تنها یک}
پردازنده اجرا شود. به بیان دیگر، کار برابر است با مجموع زمان صرف‌شده توسط هر بلوک. اگر هر بلوک یک واحد زمان طول بکشد، آنگاه کار برابر با تعداد رأس‌های گراف خواهد بود.
\itm
اما «گستره»، سریع‌ترین زمان ممکن برای اجرای محاسبه بر روی تعداد نامحدودی پردازنده است، که متناظر است با مجموع زمان بلوک‌های طولانی‌ترین مسیر در گراف توازی است.
در محاسبه «طولانی‌ترین» مسیر هر بلوک به اندازه‌ی زمان اجرای خودش وزن‌دهی می‌شود.
به این مسیر، «مسیر بحرانی» گفته می‌شود و بنابراین گستره برابر است با وزن مسیر بحرانی
\fn{critical path}
در گراف توازی.
\end{itemframe-s}


\begin{itemframe-s}{مبانی موازی‌سازی انشعاب-الحاق}{معیارهای کارایی}
\itm
برای یک گراف توازی که هر بلو‌ک آن یک واحد زمان طول می‌کشد، گستره برابر است با تعداد بلوک‌ها در مسیر بحرانی.
\itm
برای مثال در شکل زیر، گراف توازی دارای 17 رأس در کل و 8 رأس در مسیر بحرانی است؛ بنابراین اگر هر بلوک یک واحد زمان طول بکشد، کار برابر 17 واحد و گستره برابر 8 واحد زمان است.

\pic[.3]{2}
\end{itemframe-s}


\begin{itemframe-s}{مبانی موازی‌سازی انشعاب-الحاق}{معیارهای کارایی}
\itm
زمان اجرای یک محاسبه‌ی به تعداد پردازنده‌های موجود و همچنین به نحوه‌ی تخصیص بلوک‌ها توسط زمان‌بند نیز وابسته است.
\itm
برای نمایش زمان اجرای یک محاسبه‌ی وظیفه-موازی بر روی P پردازنده، از اندیس زیرنویس استفاده می‌کنیم. برای مثال، زمان اجرای یک الگوریتم روی P پردازنده را
$T_P$
می‌نویسیم.
\itm
دیدم که «کار» همان زمان اجرا روی یک پردازنده است پس برابر است با
$T_1$ .
همچنین گستره همان زمان اجرا روی تعداد نامحدودی پردازنده است؛ یعنی اگر هر بلوک روی پردازنده‌ی خودش اجرا شود. بنابراین گستره را با
$T_∞$
نشان می‌دهیم.
\end{itemframe-s}


\begin{itemframe-s}{مبانی موازی‌سازی انشعاب-الحاق}{معیارهای کارایی}
\itm
کار و گستره، کران‌های پایین زمان اجرای یک محاسبه را تعیین می‌کنند:
\itm[1]
«قانون کار»
\fn{Work Law}:
یک کامپیوتر موازی با P پردازنده در یک واحد زمان، حداکثر می‌تواند P واحد کار انجام دهد. بنابراین در زمان
$T_P$
کل کاری که می‌تواند انجام شود حداکثر
$P T_P$
است. از آنجا که کل کار برابر
$T_1$
است، خواهیم داشت:
$$P T_P ≥ T_1$$
با تقسیم بر P می‌رسیم به:

$$T_P ≥ T_1 / P$$
که به آن قانون کار گفته می‌شود.
\end{itemframe-s}


\begin{itemframe-s}{مبانی موازی‌سازی انشعاب-الحاق}{معیارهای کارایی}
\itm[2]
«قانون گستره»
\fn{span Law}:
یک کامپیوتر موازی با P
پردازنده نمی‌تواند سریع‌تر از ماشینی با تعداد نامحدود پردازنده باشد. بنابراین:
$$T_P ≥ T_∞$$

به نامساوی بالا قانون گستره می‌گوییم.
\end{itemframe-s}


\begin{itemframe-s}{مبانی موازی‌سازی انشعاب-الحاق}{معیارهای کارایی}
\itm
شتاب
\fn{speedup}
یک محاسبه روی P پردازنده را به این صورت تعریف می‌کنیم:
$$T_1 / T_P$$
که بیان می‌کند محاسبه چند برابر سریع‌تر از حالت تک‌پردازنده اجرا می‌شود.
\itm
با استفاده از قانون کار، چون
$$T_P ≥ T_1 / P$$
داریم:
$$T_1 / T_P ≤ P$$
بنابراین، شتاب روی یک کامپیوتر موازی با P
پردازنده، حداکثر برابر با P
خواهد بود.
\end{itemframe-s}



\begin{itemframe-s}{مبانی موازی‌سازی انشعاب-الحاق}{معیارهای کارایی}
\itm
می‌گوییم محاسبه «شتاب خطی»
\fn{linear speedup}
دارد. اگر
$$T_1 / T_P = \Theta(P)$$
باشد،
«شتاب خطی کامل»
\fn{perfect linear speedup}
زمانی رخ می‌دهد که:
$$T_1 / T_P = P$$
\end{itemframe-s}


\begin{itemframe-s}{مبانی موازی‌سازی انشعاب-الحاق}{معیارهای کارایی}
\itm
نسبت کار به گستره، «نرخ موازی‌سازی»
\fn{parallelism}
محاسبه‌ی موازی را نشان می‌دهد:
$$T_1 / T_∞$$
این کمیت را می‌توان از سه دیدگاه بررسی کرد:
\itm[1]
به‌عنوان یک نسبت: مقدار متوسط کاری که در هر گام مسیر بحرانی می‌تواند به‌طور موازی انجام شود.
\itm[2]
به‌عنوان یک کران بالا: بیشینه‌ی شتابی که روی هر تعداد پردازنده می‌تواند به دست آید.
\itm[3]
به‌عنوان یک محدودیت بنیادی: اگر تعداد پردازنده‌ها بیشتر از «نرخ موازی‌سازی» شود، رسیدن به شتاب خطی کامل غیرممکن خواهد بود.
\end{itemframe-s}


\begin{itemframe-s}{مبانی موازی‌سازی انشعاب-الحاق}{معیارهای کارایی}
\itm
دیدگاه سوم از بقیه مهم‌تر است؛ برای روشن شدن آن، فرض کنید:
$$P > T_1 / T_∞$$
در این صورت، قانون گستره نشان می‌دهد:
$$T_1 / T_P ≤ T_1 / T_∞ < P$$
بنابراین:
$$T_1 / T_P < P$$
بنابراین شتاب نمی‌تواند برابر با تعداد پردازنده‌ها باشد.
\end{itemframe-s}



\begin{itemframe-s}{مبانی موازی‌سازی انشعاب-الحاق}{معیارهای کارایی}
\itm
به‌علاوه، اگر تعداد پردازنده‌ها خیلی بیشتر از نرخ موازی‌سازی باشد؛ یعنی
$$P ≫ T_1 / T_∞$$
آنگاه:

$$T_1 / T_P ≪ P$$
یعنی شتاب بسیار کمتر از تعداد پردازنده‌ها خواهد بود.
\itm
به بیان دیگر، اگر تعداد پردازنده ها بیشتر از نرخ موازی‌سازی باشد،
افزودن پردازنده‌ها شتاب ما را از شتاب کامل دورتر می‌کند.
\end{itemframe-s}


\begin{itemframe-s}{مبانی موازی‌سازی انشعاب-الحاق}{معیارهای کارایی}
\itm
برای مثال بار دیگر شکلی که پیش‌تر دیدیم را در نظر بگیرید و فرض کنید هر بلوک یک واحد زمان طول می‌کشد.
\pic{2}
\end{itemframe-s}


\begin{itemframe-s}{مبانی موازی‌سازی انشعاب-الحاق}{معیارهای کارایی}
\itm
دیدم که
$T_1 = 17 $
و
$ T_∞ = 8 $
بنابراین نرخ موازی‌سازی برابر است با:

$$T_1 / T_∞ = 17 / 8 = 2.125$$
\itm
در نتیجه، فارغ از اینکه چند پردازنده استفاده شود، دست‌یابی به بیش از دو برابر کارایی (شتاب) غیرممکن است.
\itm
دقت کنید که این تنها برای
$P-FIB(4)$
بود. اما برای ورودی‌های بزرگ‌تر،
نرخ موازی‌سازی بالایی خواهیم داشت.
\end{itemframe-s}


\begin{itemframe-s}{مبانی موازی‌سازی انشعاب-الحاق}{معیارهای کارایی}
\itm
رابطه‌ی میان موازی‌سازی یک محاسبه و تعداد پردازنده‌های موجود، مفهوم مهمی به نام انعطاف‌پذیری
\fn{slackness}
را معرفی می‌کند.
\itm
انعطاف‌پذیری یک محاسبه روی تعداد P
پردازنده، برابر است با:
$$\frac{T_1 / T_\infty}{P} \;=\; \frac{T_1}{P T_\infty}$$
که نشان می‌دهد موازی‌سازی محاسبه تا چه اندازه از تعداد پردازنده‌های موجود در ماشین بیشتر است.

\end{itemframe-s}


\begin{itemframe-s}{مبانی موازی‌سازی انشعاب-الحاق}{معیارهای کارایی}
\itm
اگر مقدار انعطاف‌پذیری کمتر از 1 باشد، شتاب خطی کامل غیرممکن است، زیرا
$$\frac{T_1}{P T_\infty} < 1$$
و طبق قانون گستره:

$$\frac{T_1}{T_P} \;\leq\; \frac{T_1}{T_\infty} \;<\; P.$$
\itm
در واقع، هرچه مقدار انعطاف‌پذیری از 1 کمتر شود و به 0 نزدیک شود، شتاب بیشتر و بیشتر از شتاب خطی کامل فاصله می‌گیرد.
در این حالت، افزودن موازی‌سازی به الگوریتم می‌تواند تأثیر زیادی بر کارایی آن داشته باشد.
\end{itemframe-s}


\begin{itemframe-s}{مبانی موازی‌سازی انشعاب-الحاق}{معیارهای کارایی}
\itm
اگر مقدار انعطاف‌پذیری بیشتر از 1 باشد، در این صورت «کار به‌ازای هر پردازنده» به محدودکننده‌ی اصلی تبدیل می‌شود.
در این حالت، یک زمان‌بند خوب می‌تواند شتابی نزدیک‌تر و نزدیک‌تر به شتاب خطی کامل نتیجه دهد.
\itm
وقتی مقدار انعطاف‌پذیری خیلی بیشتر از ۱ شود، اضافه کردن پردازنده‌های بیشتر، دیگر فایده‌ی چندانی ندارد و سود ناشی از موازی‌سازی به‌تدریج کمتر و کمتر می‌شود.
\end{itemframe-s}


\begin{itemframe-s}{مبانی موازی‌سازی انشعاب-الحاق}{زمان‌بندی}
\itm
عملکرد خوب تنها به کمینه‌سازی کار و گستره وابسته نیست.
باید بلوک‌ها نیز به‌طور کارآمد روی پردازنده‌های ماشین موازی زمان‌بندی شوند.
\itm
در مدل برنامه‌نویسی انشعاب-الحاق نمی‌توان مشخص کرد کدام بلوک روی کدام پردازنده اجرا شود.
در عوض، ما به زمان‌بند تکیه می‌کنیم تا محاسبات را به صورت پویا را به پردازنده‌ها نگاشت کند.
\itm
در عمل، زمان‌بند بلوک‌ها را به نخ‌ها نگاشت می‌کند و سپس سیستم‌عامل این نخ‌ها را روی پردازنده‌ها زمان‌بندی می‌کند.
اما ما می‌توانیم به‌سادگی تصور کنیم که زمان‌بند بلوک‌ها را مستقیماً به پردازنده‌ها نگاشت می‌کند.
\itm
یک زمان‌بند باید «در لحظه»
\fn{online}
عمل کند. یعنی بدون اینکه از قبل بداند
چه زمانی فراخونی‌های موازی (spawn) ایجاد می‌شوند یا چه زمانی پایان می‌یابند، محاسبه را زمان‌بندی کند.
\end{itemframe-s}


\begin{itemframe-s}{مبانی موازی‌سازی انشعاب-الحاق}{زمان‌بندی}
\itm
به‌طور خاص، ما زمان‌بندهای حریصانه
\fn{greedy schedulers}
را تحلیل می‌کنیم که در هر گام بیشترین تعداد بلوک را به پردازنده‌ها اختصاص می‌دهند و سعی می‌کنند هیچ پردازنده‌ای را بیکار نگذارند.
در یک زمان‌بند حریصانه برای هر گام دو حالت وجود دارد:
\itm[1]
 گام کامل
\fn{Complete step}
: دست‌کم $P$ بلوک, آماده‌ی اجرا هستند. یک زمان‌بند حریصانه هر $P$ بلوک آماده را به پردازنده‌ها تخصیص می‌دهد.
\itm[2]

گام ناقص
\fn{Incomplete step}
: کمتر از $P$ بلوک آماده‌ی اجرا هستند. زمان‌بند حریصانه هر بلوک‌ آماده را به پردازنده‌ی خودش تخصیص می‌دهد، تعدادی پردازنده بیکار باقی می‌مانند، اما همه‌ی بلوک‌های آماده اجرا می‌شوند.
\itm
منظور از بلوک آماده بلوکی است که همه بلوک‌هایی که به آنها وابسته است اجرا شده باشند.
\end{itemframe-s}


\begin{itemframe-s}{مبانی موازی‌سازی انشعاب-الحاق}{زمان‌بندی}
\itm
دیدیم که طبق قانون کار، سریع‌ترین زمان اجرای یک محاسبه روی $P$ پردازنده دست‌کم برابر است با:
$$T_P ≥ T_1 / P$$
از طرف دیگر، قانون گستره بیان می‌کند که سریع‌ترین زمان ممکن دست‌کم برابر است با:
$$T_P ≥ T_∞$$
\itm

در ادامه قضیه‌ایی را بررسی می‌کنیم که نشان می‌دهد زمان اجرای یک محاسبه با زمان‌بند حریصانه کران بالایی برابر با مجموع این دو کران پایین دارد.
این کران بالا خوب است زیرا تا حد قابل قبولی به کران‌های پایین نزدیک است.
\end{itemframe-s}


\begin{itemframe-s}{مبانی موازی‌سازی انشعاب-الحاق}{زمان‌بندی}
\itm
می‌خواهیم ثابت کنیم زمان اجرای یک محاسبه با زمانبند حریصانه حداکثر به این مقدار زمان نیاز دارد:
$$T_P ≤ T_1 / P + T_∞$$
\itm
\textbf{اثبات}:
فرض کنید هر بلوک یک واحد زمان طول می‌کشد. (هر بلوک طولانی‌تر را می‌توان با زنجیره‌ای از بلوک‌های تک‌واحدی جایگزین کرد.) حال گام‌های کامل و ناقص را جداگانه بررسی می‌کنیم:
\itm[الف]
در هر گام کامل، $P$ پردازنده روی‌هم $P$ واحد کار انجام می‌دهند. بنابراین، اگر تعداد گام‌های کامل $k$ باشد، کل کاری که در این گام‌ها انجام می‌شود $kP$ است. این مقدار از تعداد کل کارها کمتر است پس
$kP ≤ T_1$
از این‌جا نتیجه می‌گیریم که تعداد گام‌های کامل حداکثر
$$T_1 / P$$
است.
\end{itemframe-s}


\begin{itemframe-s}{مبانی موازی‌سازی انشعاب-الحاق}{زمان‌بندی}
\itm[ب]
حال گام ناقص را در نظر بگیریم. فرض کنید $G$ گراف توازی کل محاسبه باشد.
در ابتدای گام ناقص، زیرگراف توازی
$G'$
هنوز اجرا نشده است و پس از این گام زیرگراف توازی $G''$ باقی می‌ماند.
\sub
همچنین تعریف می‌کنیم مجموعهُ $R$ هم بلوک‌هایی هستند که در آغاز این گام ناقص آماده‌اند.
\sub
یک طولانی‌ترین مسیر در
$G'$
باید با بلوک‌ای از
$R$
شروع شود. فرض کنید یک طولانی‌ترین مسیر در
$G'$
از بلوکی خارج از R شروع شود، می‌دانیم این بلوک آماده نیست و گراف بدون دور است پس می‌توانیم پیش‌نیاز این بلوک را به مسیر اضافه کنیم و در اینحا به تناقض می‌رسیم.

\end{itemframe-s}


\begin{itemframe-s}{مبانی موازی‌سازی انشعاب-الحاق}{زمان‌بندی}
\itm
چون زمان‌بند حریصانه همه‌ی بلوک‌های آماده را در این گام اجرا می‌کند،
$G''$
برابر است با
$G'$
 منهای
$R$.
\itm

بنابراین طولانی‌ترین مسیر در
$G''$
دقیقاً یک واحد کوتاه‌تر از طولانی‌ترین مسیر در
$G'$
است. در نتیجه، هر گام ناقص طول مسیر بحرانی زیر گراف باقی‌مانده را یک واحد کاهش می‌دهد.
\itm
می‌دانیم که طول مسیر بحرانی یا همان گستره برابر است با
$T_∞$
بنابراین تعداد گام‌های ناقص حداکثر
$T_∞$
است.
\end{itemframe-s}


\begin{itemframe-s}{مبانی موازی‌سازی انشعاب-الحاق}{زمان‌بندی}
\itm
دیدیم که حداکثر تعداد گام‌های کامل $T_1 / P$ است و حداکثر تعداد گام‌های ناقص $T_∞$ است بنابراین حداکثر تعداد مجموع گام‌ها
$$T_∞+T_1 / P$$
است. و بدین ترتیب قضیه ثابت می‌شود.
\itm
با استفاده از این قضیه ثابت می‌شود که اگر
$$P \gg T_1 / T_∞$$
باشد، (یعنی انعطاف‌پذیری خیلی بیشتر از ۱ باشد) آنگاه،
شتاب تقریباً برابر با $P$ است (شتاب نزدیک به کامل).
\end{itemframe-s}


\begin{itemframe-s}{مبانی موازی‌سازی انشعاب-الحاق}{تحلیل الگوریتم‌های موازی}
\itm
اکنون همه‌ی ابزارهای لازم را برای تحلیل الگوریتم‌های موازی با استفاده از تحلیل کار-گستره را در اختیار داریم.
\itm
تحلیل «کار» ساده است. کافیست تصویر سریال الگوریتم موازی را به دست آورده و با همان روش تحلیل زمانی الگوریتم‌های سری آن را تحلیل کنیم.

\itm
آنچه جدید است، تحلیل گستره است. اما خواهید دید که چندان پیچیدگی بیشتری ندارد.
\end{itemframe-s}


\begin{itemframe-s}{مبانی موازی‌سازی انشعاب-الحاق}{تحلیل الگوریتم‌های موازی}
\itm
بیایید دوباره به مثال P-FIB بازگردیم و با این مثال با بعضی از ایده‌های ابتدایی آشنا شویم.
\itm
دیدم که برای تحلیل کار باید تصویر سری الگوریتم P-FIB را به دست آورده و آن را تحلیل زمانی عادی کنیم. تصویر سری این الگوریتم همان الگوریتم FIB است که پیش‌تر آن را تحیلی کرده‌ایم بنابراین می‌دانیم:
$$T_1(n) = T(n) = \Theta(\phi^n)$$
\end{itemframe-s}


\begin{itemframe-s}{مبانی موازی‌سازی انشعاب-الحاق}{تحلیل الگوریتم‌های موازی}
\itm
حال با استفاده از شکل زیر نشان خواهیم داد که چگونه باید گستره را تحلیل کنیم.
\pic[.8]{3}
\end{itemframe-s}


\begin{itemframe-s}{مبانی موازی‌سازی انشعاب-الحاق}{تحلیل الگوریتم‌های موازی}
\itm[1]
اگر دو محاسبه به‌صورت سری ترکیب شوند، گستره‌ی کل برابر است با جمع گستره دو محاسبه.
\itm[2]
اما اگر به‌صورت موازی ترکیب شوند، گستره کل برابر است با بیشینه‌ی گستره‌های دو محاسبه.
\itm
در واقع، گراف توازی هر محاسبه‌ای را می‌توان با ترکیب موازی یا سری تعدادی بلوک ساخت.
\itm
حال که ترکیب سری و موازی را آموختیم، می‌توانیم گستره‌ی
$P-FIB(n)$
را تحلیل کنیم.
\end{itemframe-s}


\begin{itemframe-s}{مبانی موازی‌سازی انشعاب-الحاق}{تحلیل الگوریتم‌های موازی}
\itm
فراخوان
$P-FIB(n-1)$
به صورت موازی با فراخوانی
$P-FIB(n-2)$
اجرا می‌شود. بنابراین، گستره‌ی
$P-FIB(n)$
را می‌توان به‌صورت رابطه‌ی بازگشتی زیر نوشت:
\begin{align*}
T_∞(n) &= \max\{T_∞(n-1), T_∞(n-2)\} + \Theta(1)\\
        &= T_∞(n-1) + \Theta(1),
\end{align*}
\itm
که جواب آن برابر است با:
$$T_∞(n) = \Theta(n)$$
\end{itemframe-s}


\begin{itemframe-s}{مبانی موازی‌سازی انشعاب-الحاق}{تحلیل الگوریتم‌های موازی}
\itm
پس نرخ موازی‌سازی در P-FIB(n) برابر است با:
$$
T_1(n) / T_∞(n) = \Omega(\phi^n / n)
$$

که با افزایش n به ‌سرعت رشد می‌کند.
\end{itemframe-s}


\begin{itemframe-s}{مبانی موازی‌سازی انشعاب-الحاق}{حلقه‌های موازی}
\itm
بسیاری از الگوریتم‌ها شامل حلقه‌هایی هستند که تمام تکرارهای آن‌ها می‌توانند به صورت موازی اجرا شوند.
\itm
هرچند می‌توان از کلیدواژه‌های spawn و sync برای موازی‌سازی چنین حلقه‌هایی استفاده کرد، ولی راحت‌تر است که کلید واژه‌ی جداگانه‌ایی داشته باشیم که مشخص کند یک حلقه می‌تواند موازی اجرا شود.
این قابلیت از طریق کلیدواژه‌ی parallel فراهم می‌شود که قبل از کلیدواژه‌ی for قرار می‌گیرد.



\end{itemframe-s}


\begin{itemframe-s}{مبانی موازی‌سازی انشعاب-الحاق}{حلقه‌های موازی}
\itm
به‌عنوان مثال، رویه‌ی
\texttt{P-MAT-VEC}
ضرب یک ماتریس مربعی در بردار را به صورت موازی انجام می‌دهد. کلیدواژه‌های parallel قبل از for نشان می‌دهند که هر تکرار از بدنه‌ی حلقه، می‌توانند به صورت موازی اجرا شوند.
\begin{algorithm}[H]\alglr
\caption{P-MAT-VEC}
\begin{algorithmic}[1]
\Func{P-MAT-VEC}{$A, x, y, n$}
  \pFor{$i = 1 \ \textbf{to} \ n$} \Comment{parallel loop}
    \For{$j = 1 \ \textbf{to} \ n$} \Comment{serial loop}
      \State $y_i \gets y_i + a_{ij} \, x_j$
    \EndFor
  \EndpFor
\end{algorithmic}
\end{algorithm}
\end{itemframe-s}


\begin{itemframe-s}{مبانی موازی‌سازی انشعاب-الحاق}{حلقه‌های موازی}
\itm
کامپایلرها برای برنامه‌های موازی fork-join می‌توانند حلقه‌های parallel for را با استفاده از spawn و sync پیاده‌سازی کنند. اینکار به شکل بازگشتی انجام می‌شود.
\itm
به‌عنوان مثال، برای حلقه‌ی موازی بالا، یک کامپایلر می‌تواند تابع کمکی
\texttt{P-MAT-VEC-RECURSIVE}
را تولید کرده و سپس حلقه for را با فراخوانی
$$\texttt{P-MAT-VEC-RECURSIVE}(A, x, y; n, ,1, n)$$

جایگزین کند.
\end{itemframe-s}


\begin{itemframe-s}{مبانی موازی‌سازی انشعاب-الحاق}{حلقه‌های موازی}
\itm
این روال بازگشتی نصف اول تکرارهای حلقه را با نصف دوم آن، موازی اجرا کرده و سپس یک sync اجرا می‌کند.
بدین ترتیب یک درخت دودویی از اجرای موازی ساخته می‌شود. خطوط 2–3 نیز حالت پایه می‌باشد.
\begin{algorithm}[H]\alglr
\caption{P-MAT-VEC-RECURSIVE}
\begin{algorithmic}[1]
\Func{P-MAT-VEC-RECURSIVE}{$A, x, y, n, i, i'$}
\If{$i = i'$} \Comment{just one iteration to do}
    \For{$j = 1 \ \textbf{to} \ n$} \Comment{mimic P-MAT-VEC serial loop}
        \State $y_i \gets y_i + a_{ij} \, x_j$
    \EndFor
\Else
    \State $mid \gets \lfloor (i + i')/2 \rfloor$ \Comment{parallel divide-and-conquer}
    \State \textbf{spawn} P-MAT-VEC-RECURSIVE$(A, x, y, n, i, mid)$
    \State P-MAT-VEC-RECURSIVE$(A, x, y, n, mid+1, i_0)$
    \State \textbf{sync}
\EndIf
\end{algorithmic}
\end{algorithm}

\end{itemframe-s}


\begin{itemframe-s}{مبانی موازی‌سازی انشعاب-الحاق}{حلقه‌های موازی}
\itm
برای محاسبه‌ی کار
\texttt{P-MAT-VEC}
کافی است زمان اجرای نسخه‌ی ترتیبی آن را حساب کنیم؛ یعنی حالتی که حلقه‌ی parallel for در خط 1 با یک حلقه‌ی معمولی for جایگزین شود. بنابراین داریم:
$$
T_1(n) = \Theta(n^2)
$$
\itm
این تحلیل به‌ظاهر سربار ناشی از فراخوانی بازگشتی برای پیاده‌سازی حلقه‌های موازی را نادیده می‌گیرد.
درست است که این فراخوانی‌ها سربار دارد، اما از نظر مجانبی تاثیری روی زمان اجرا نمی‌گذارد.
\end{itemframe-s}


\begin{itemframe-s}{مبانی موازی‌سازی انشعاب-الحاق}{حلقه‌های موازی}
\itm
در هنگام تحلیل کار دیدم که می‌توانیم از سربار فراخوانی‌های بازگشتی صرف نظر کنیم. ولی هنگام تحلیل «گستره» باید آن را لحاظ کرد.
\itm
به طور کلی اگر یک حلقه موازی با n تکرار داشته باشیم و گستره تکرار iام آن را
$iter_\infty(i)$
بنامیم، گستره حلقه برابر خواهد بود با:
$$
T_\infty(n) = \Theta(\lg n) + \max \{ iter_\infty(i) \mid 1 \leq i \leq n \}
$$
\itm
برای فهم بهتر تساوی بالا، توجه کنید که گه عمق درخت بازگشت لگاریمی است. حال باید تاثیر عملیاتی که درون حلقه انجام می‌شود را لحاظ کنیم که همان عبارت دوم است. در حالت سری ما تاثیر حلقه دوم را در تکرار حلقه اول ضرب می‌کردیم. در اینجا دقت کنید که به خاطر موازی بودن تکرارهای حلقه، جمع صورت گرفته.
\end{itemframe-s}


\begin{itemframe-s}{مبانی موازی‌سازی انشعاب-الحاق}{حلقه‌های موازی}
\itm
حال که تحلیل حلقه‌های موازی را در حالت کلی آموختیم بیایید به تحیلیل الگوریتم
\texttt{P-MAT-VEC}
بازگردیم.
\itm

گستره حلقه‌ی بیرونی بدون در نظر گرفتن بدنه حلقه برابر است با:
$$\Theta(\lg n)$$
\itm
 در هر تکرار از حلقه‌ی بیرونی، حلقه‌ی درونی شامل n تکرار است و در همه تکرارهای حلقه‌ی بیرونی، گستره حلقه درونی برابر است.
پس نتیجه گرفتن بیشینهٔ گستره حلقه درونی، روی تمام تکرارهای حلقه‌ی بیرونی برابر است با:
$$\Theta(n)$$

\end{itemframe-s}


\begin{itemframe-s}{مبانی موازی‌سازی انشعاب-الحاق}{حلقه‌های موازی}
\itm
بنابراین، گستره‌ی کلی این تابع برابر است با:
$$T_\infty(n) = \Theta(n) + \Theta(\lg n) = \Theta(n)$$
\itm
دیدم که کار برابر
$Θ(n^2)$
است پس نرخ موازی‌سازی برابر خواهد بود با:
$$
\Theta(n^2) / \Theta(n) = \Theta(n).
$$
\end{itemframe-s}


\begin{itemframe-s}{مبانی موازی‌سازی انشعاب-الحاق}{وضعیت رقابت}
\itm
یک الگوریتم موازی را قطعی
\fn{deterministic}
می‌نامیم اگر همواره روی یک ورودی مشخص، یک خروجی یکسان تولید کند.
اگر خروجی الگوریتم بتواند بین اجراهای مختلف روی یک ورودی تغییر کند، الگوریتم غیرقطعی
\fn{nondeterministic}
است.
\itm
زمانی که یک الگوریتم را به صورت موازی پیاده‌سازی می‌کنیم ممکن است خطایی به نام «عدم قطعیت رقابتی»
\fn{determinacy race}
رخ دهد که یکی از انواع «وضعیت رقابت»
\fn{Race condition}
است.
\itm
خطای عدم قطعیت رقابتی می‌تواند باعث شود یک الگوریتم موازی که هدف آن قطعی‌بودن است، به‌طور غیرقطعی عمل کند.
تشخیص این خطا می‌تواند بسیار سخت باشد.
\end{itemframe-s}


\begin{itemframe-s}{مبانی موازی‌سازی انشعاب-الحاق}{وضعیت رقابت}
\itm
خطاهای مشهور زیر در طول تاریخ به از نوع «عدم قطعیت رقابتی» بوده اند:
\itm[1]
خطای دستگاه پرتودرمانی Therac-25 که سه نفر را کشت و چندین نفر دیگر را مجروح کرد.
\itm[2]
خاموشی بزرگ شمال‌شرق آمریکا در سال 2003 که باعث قطعی برق بیش از 50 میلیون نفر شد.
%(البته این که برای ما چیز عجیبی نیست)
\itm
ممکن است نرم‌افزار شما روزها در آزمایشگاه بدون خطا اجرا شود، ولی سپس به‌طور پراکنده در محیط واقعی دچار کرش‌های جدی گردد.

\end{itemframe-s}


\begin{itemframe-s}{مبانی موازی‌سازی انشعاب-الحاق}{وضعیت رقابت}
\itm

عدم قطعیت رقابتی هنگامی رخ می‌دهد که دو دستور موازی به یک مکان حافظه‌ی مشترک دسترسی یابند و دست‌کم یکی از آن‌ها مقدار ذخیره‌شده را تغییر دهد. تابع ساده‌ی زیر این وضعیت را نشان می‌دهد.
\begin{algorithm}[H]\alglr
\caption{RACE-EXAMPLE}
\begin{algorithmic}[1]
\Func{RACE-EXAMPLE}{}
\State $x \gets 0$
\pFor{$i = 1 \ \textbf{to} \ 2$}
    \State $x \gets x + 1$ \Comment{determinacy race}
\EndpFor
\State print($x$)
\end{algorithmic}
\end{algorithm}
\end{itemframe-s}


\begin{itemframe-s}{مبانی موازی‌سازی انشعاب-الحاق}{وضعیت رقابت}
\itm
تابع
\texttt{RACE-EXAMPLE}
دو بلوک موازی ایجاد می‌کند که هر یک مقدار x را یک واحد افزایش می‌دهند.
هرچند انتظار داریم این روال مانند نسخه ترتیبی، همواره مقدار 2 را چاپ کند ولی ممکن است مقدار 1 چاپ شود.
زیرا عمل افزایش مقدار x در کد سطح پایین‌تر (اسمبلی) در واقع از چند دستور تشکیل شده است:
\itm[1]
بارگذاری x
از حافظه به یکی از ثبات‌های پردازنده.
\itm[2]
افزایش مقدار در ثبات.
\itm[3]
ذخیره‌سازی مقدار جدید ثبات در حافظه.
\itm
اگر دستورات دو پردازنده به تریتیب خاصی «درهم‌تنیده»
\fn{interleaved}
شوند، تاثیر یکی از به‌روزرسانی‌ها ختنثی می‌شود و در نهایت مقدار 1 چاپ خواهد شد.
\end{itemframe-s}


\begin{itemframe-s}{مبانی موازی‌سازی انشعاب-الحاق}{وضعیت رقابت}
\itm
توجه کنید که حتی با وجود خطای عدم قطعیت رقابتی، باز هم  بسیاری از اجراها منجر به نتیجه‌ی درست می‌شوند.
اما برخی اجراها به‌طور غیرقابل پیش‌بینی نتیجه‌ی غلط تولید می‌کنند.
همین غیرقابل‌تکرار بودن، یافتن خطاهای رقابتی را به‌شدت دشوار می‌سازد.
\end{itemframe-s}


\begin{itemframe-s}{مبانی موازی‌سازی انشعاب-الحاق}{وضعیت رقابت}
\itm
شکل زیر گراف توازی الگوریتم
\texttt{RACE-EXAMPLE}
را نمایش می‌دهد، با این تفاوت که بلوک‌ها تا سطح دستورات سطح پایین شکسته شده‌اند.
\pic[.4]{4}
\end{itemframe-s}


\begin{itemframe-s}{مبانی موازی‌سازی انشعاب-الحاق}{وضعیت رقابت}
\itm
گفتیم که یک طبق فرض این بحث، کامپیوتر ما از سازگاری ترتیبی
\fn{sequential consistency}
پشتیبانی می‌کند. بنابراین می‌توانید اجرای موازی دو بلوک را مانند یک توالی خطی از دستورات هر دو بلوک درنظر بگیرید.
\itm
به این توالی خطی درهم‌تنیدگی
\fn{interleaving}
گفته می‌شود زیرا چند دستور موازی می‌توانند در هم تنیده شوند بدین معنی که دستورات سطح پایین می‌توانند در میان یکدیگر اجرا شوند.
البته به‌گونه‌ای که وابستگی‌های موجود در گراف توازی حفظ شوند.

\end{itemframe-s}


\begin{itemframe-s}{مبانی موازی‌سازی انشعاب-الحاق}{وضعیت رقابت}
\itm
شکل زیر، یک ترتیب خاص از اجرای دستورات را نشان می‌دهد که باعث ایجاد این مشکل می‌شود.
\pic[.4]{5}
\end{itemframe-s}


\begin{itemframe-s}{مبانی موازی‌سازی انشعاب-الحاق}{وضعیت رقابت}
\itm
هر چند با روش‌هایی مانند انحصار متقابل
\fn{mutual-exclusion locks}
، می‌توان از این خطا جلوگیری کرد. اما ما در این بحث به طور کلی وجود وضعیت رقابت را مجاز نمی‌دانیم و می‌خواهیم ببینم چطور می‌توان از بروز این وضعیت جلوگیری کرد.
\itm
برای اطمینان از قطعی بودن الگوریتم‌ها، هر دو رشته‌ای که به صورت موازی عمل می‌کنند باید بدون تداخل متقابل
\fn{mutually noninterfering}
باشند. یعنی هر دو تنها بخوانند و هیچ مقدار مشترکی را تغییر ندهند.
\itm
بنابراین در یک ساختار parallel for باید تمام تکرارهای بدنه حلقه، بدون تداخل متقابل باشند. همچنین بین یک spawn و sync متناظر، نباید تداخل متقابل وجود داشته باشد.

\end{itemframe-s}


\begin{itemframe-s}{مبانی موازی‌سازی انشعاب-الحاق}{وضعیت رقابت}
\itm
بیایید یک مثال ساده از پیاده‌سازی اشتباه تابع
\texttt{P-MAT-VEC}
ببینیم.
تابع زیر را در نظر بگیرید. این روال حلقه‌ی درونی را نیز موازی کرده و گستره‌ٔ
$\Theta(lgn)$
را به‌دست می‌آورد، اما به دلیل وضعیت رقابت در خط ۳، خروجی نادرست تولید می‌کند:
\begin{algorithm}[H]\alglr
\caption{P-MAT-VEC-WRONG}
\begin{algorithmic}[1]
\Func{P-MAT-VEC-WRONG}{$A, x, y, n$}
\pFor{$i = 1 \ \textbf{to} \ n$}
    \pFor{$j = 1 \ \textbf{to} \ n$}
        \State $y_i \gets y_i + a_{ij} \, x_j$ \Comment{determinacy race}
    \EndpFor
\EndpFor
\end{algorithmic}
\end{algorithm}
\itm
توجه کنید در اینجا متغیرهای اندیس
\fn{index}
مانند i و j مشکلی ایجاد نمی‌کنند. زیرا هر تکرار عملاً نسخه‌ی مستقل خودش از متغیر اندیس را ایجاد می‌کند.
\end{itemframe-s}


\begin{itemframe-s}{مبانی موازی‌سازی انشعاب-الحاق}{مثالی از تحیلی کار-گستره}
\itm
رای نشان دادن قدرت تحلیل کار-گستره، یک داستان واقعی را بررسی خواهیم کرد که چندین سال پیش در جریان توسعه یکی از اولین برنامه‌های شطرنج موازی سطح جهانی رخ داده است.
\itm
این برنامه شطرنج روی یک رایانه ۳۲ هسته‌ایی توسعه و آزمایش شد، اما قرار بود روی یک ابررایانه با ۵۱۲ هسته اجرا شود.
از آنجا که دسترسی به ابررایانه محدود و پرهزینه بود، توسعه‌دهندگان آزمایش‌ها را روی رایانه کوچک اجرا کرده و عملکرد را برای رایانه بزرگ برآورد می‌کردند.
\itm
در مقطعی، توسعه‌دهندگان یک بهینه‌سازی به برنامه اضافه کردند که زمان اجرای برنامه را روی رایانه کوچک از
$T_{32} = 65$
ثانیه به
$T'_{32} = 40$
کاهش داد. با این حال، با استفاده از تحلیل کار و گستره، نتیجه گرفتند که نسخه بهینه‌شده که روی ۳۲ پردازنده سریع‌تر بود، روی ۵۱۲ پردازنده از نسخه اصلی کندتر خواهد بود. بنابراین، آنها از اجرای این بهینه‌سازی صرف‌نظر کردند.
\end{itemframe-s}


\begin{itemframe-s}{مبانی موازی‌سازی انشعاب-الحاق}{مثالی از تحیلی کار-گستره}
\itm
تحلیل کار-گستره آنها به این صورت بود: نسخه اصلی برنامه دارای
$T_1 = 2048 $
و
$T_\infty = 1$
ثانیه بود.
\itm
نابرابری زیر را قبلاً در هنگام معرفی زمانبند دیده‌ایم:

$$T_P \le \frac{T_1}{P} + T_\infty$$
\itm
اگر این  نابرابری را به صورت برابری در نظر بگیریم می‌توانیم با جایگذاری کار و گستره، زمان اجرا روی P پردازنده را تخمین بزنیم(در واقع یک کران برای آن می‌یابیم)

\end{itemframe-s}


\begin{itemframe-s}{مبانی موازی‌سازی انشعاب-الحاق}{مثالی از تحیلی کار-گستره}
\itm
پس بیایید زمان حالت اولیه روی کامپیوتر ۳۲ هسته‌ایی را تخمین بزنیم:
$$
T_{32} = \frac{2048}{32} + 1 = 65
$$
\itm
حال به انجام بهینه‌سازی، کار به
$T'_1 = 1024$
و گستره به
$T'_\infty = 8$
ثانیه تغییر می‌یابد. پس داریم:
$$
T'_{32} = \frac{1024}{32} + 8 = 40
$$
\end{itemframe-s}



\begin{itemframe-s}{مبانی موازی‌سازی انشعاب-الحاق}{مثالی از تحیلی کار-گستره}
\itm
بنابراین روی کامپیوتر ۳۲ هسته‌ایی بهبود حاصل شد. اما نسبت سرعت دو نسخه هنگامی که زمان اجرای آنها روی ۵۱۲ پردازنده برآورد می‌شود، معکوس می‌شود. نسخه اول دارای زمان اجرای
$$
T_{512} = \frac{2048}{512} + 1 = 5
$$
و نسخه دوم دارای زمان اجرای
$$
T'_{512} = \frac{1024}{512} + 8 = 10
$$
ثانیه است.
\end{itemframe-s}


\begin{itemframe-s}{مبانی موازی‌سازی انشعاب-الحاق}{مثالی از تحیلی کار-گستره}
\itm
بهینه‌سازی‌ایی که برنامه را روی ۳۲ پردازنده سریع‌تر می‌کرد، باعث شد برنامه روی ۵۱۲ پردازنده دو برابر طولانی‌تر اجرا شود!
\itm
گستره نسخه بهینه‌شده برابر با ۸ بود که در زمان اجرای روی ۳۲ پردازنده نقش غالب نداشت، اما روی ۵۱۲ پردازنده نقش غالب پیدا می‌کند و مزیت استفاده از پردازنده‌های بیشتر را از بین می‌برد. اصطلاحاً می‌گوییم این بهینه‌سازی مقیاس‌پذیری نیست.
\itm
بنابراین در هنگام برآورد مقیاس‌پذیری، تحلیل کار-گستره نسبت به اندازه‌گیری زمان اجرا برتری دارد.
\end{itemframe-s}